\documentclass[10pt]{article}
\usepackage[a4paper,margin=23mm]{geometry}
\usepackage{amsmath,amssymb,amsthm}

\newtheorem{lemma}{Lemma}
\newtheorem{theorem}[lemma]{Theorem}
\newcommand{\Z}{\mathbb Z}
\newcommand{\Q}{\mathbb Q}
\newcommand{\U}{\operatorname U}
\newcommand{\dlt}{\delta}
\newcommand{\gm}{\gamma}
\newcommand{\Hom}{\operatorname {Hom}}

\title{The Plotkin equality for finitely generated Lie rings}
\author{}
\date{}

\begin{document}
\maketitle

Let $\U(L)$ be the integral enveloping algebra of a Lie ring $L$, and let
$\mathfrak d(L)$ be its augmentation ideal.  Since the canonical map
$L\to\U(L)$ need not be injective, define
\[
 \dlt_n(L)=\{x\in L:\text{the image of $x$ in $\U(L)$ lies in
                         $\mathfrak d(L)^n$}\}.
\]
Also set
\[
 \gm_1(L)=L,\qquad \gm_{n+1}(L)=[\gm_n(L),L],
\]
and
\[
 \dlt_\omega(L)=\bigcap_{n\ge1}\dlt_n(L),\qquad
 \gm_\omega(L)=\bigcap_{n\ge1}\gm_n(L).
\]
Throughout, ``finitely generated'' means finitely generated as a Lie ring.

\begin{theorem}\label{thm:main}
If $L$ is a finitely generated Lie ring, then for every $n\ge1$ there is an
integer $m=m(L,n)$ such that
\[
                         \dlt_m(L)\subseteq\gm_n(L).   \tag{1}\label{eq:main}
\]
Consequently $\dlt_\omega(L)=\gm_\omega(L)$.
\end{theorem}

The proof will be assembled from short lemmas.  The construction below uses
only additive decompositions, enveloping algebras, and matrices.

\begin{lemma}[central separation]\label{lem:separate}
Let $K$ be a finitely generated nilpotent Lie ring and let
$A\cong\Z/p$ be central in $K$, where $p$ is prime.  There are a nilpotent
associative ring $R$, not necessarily with an identity, and a Lie
homomorphism $\rho:K\to R^{(-)}$ whose restriction to $A$ is injective.
Here $R^{(-)}$ denotes $R$ with bracket $[a,b]=ab-ba$.
\end{lemma}

We prove Lemma~\ref{lem:separate} after the following short preparations.

\begin{lemma}[additive finite generation]\label{lem:additive-finite}
Every finitely generated nilpotent Lie ring is finitely generated as an
abelian group.
\end{lemma}

\begin{proof}
If the ring is generated by $x_1,\ldots,x_d$ and has class $c$, its additive
group is spanned by bracket words in the $x_i$.  Words of length greater
than $c$ lie in $\gm_{c+1}$ and vanish, and there are only finitely many
words of lengths at most $c$.
\end{proof}

An element of $\Hom_{\Z}(\U(H),D)$ is simply an additive function
$\U(H)\to D$.  The next elementary lemma records the useful extension
property of this group of functions.

\begin{lemma}[extending maps]\label{lem:coinduced-extension}
Let $H$ be a Lie ring, let $D$ be an abelian group, and set
\[
 J=\Hom_{\Z}(\U(H),D),\qquad (x\alpha)(u)=\alpha(ux).
\]
Let $M'\subseteq M$ be $\U(H)$-modules.  If the inclusion has an additive
retraction $M\to M'$, then every $\U(H)$-module map $M'\to J$ extends to a
$\U(H)$-module map $M\to J$.
\end{lemma}

\begin{proof}
The formula for the action is valid because
$(x(y\alpha)-y(x\alpha))(u)=\alpha(u(xy-yx))=\alpha(u[x,y])$.
Evaluation at $1$ gives a bijection
\[
 \Hom_{\U(H)}(M,J)\cong\Hom_{\Z}(M,D).
\]
Its inverse sends an additive map $h:M\to D$ to the module map
$m\mapsto(u\mapsto h(um))$.  Now start with a module map $M'\to J$,
translate it to an additive map $M'\to D$, compose with the given retraction
$M\to M'$, and translate back.  The resulting module map $M\to J$ is the
required extension.
\end{proof}

We next enlarge the central subgroup by a finite amount.  The only purpose
of the enlargement is to make the resulting additive extension split.

\begin{lemma}[splitting the additive extension]\label{lem:finite-pushout}
Let $A=\langle a\rangle\cong\Z/p$ be central in a Lie ring $K$, and suppose
that the additive group of $H=K/A$ is finitely generated.  There are a finite
cyclic $p$-group $D$ and a Lie ring $E$ containing $K$ such that $D$ is
central, $E/D\cong H$, and $E=D\oplus H$ as an abelian group.
\end{lemma}

\begin{proof}
Write the additive group of $H$ as
\[
 \Z^b\oplus\Z/n_1\oplus\cdots\oplus\Z/n_s.
\]
Choose $r$ so large that the $p$-part of every $n_i$ divides $p^{r-1}$,
and let $D=\langle d\rangle\cong\Z/p^r$.  Identify $A$ with the subgroup
generated by $p^{r-1}d$, and define
\[
 E=(D\oplus K)/\langle(p^{r-1}d,-a)\rangle .            \tag{2}\label{eq:E}
\]
The bracket $[(d_1,k_1),(d_2,k_2)]=(0,[k_1,k_2])$ descends to $E$ because
$A$ is central.  The natural maps from $D$ and $K$ to $E$ are injective:
a multiple of $(p^{r-1}d,-a)$ can have either coordinate zero only when
that multiple is divisible by $p$, in which case both coordinates are zero.
Clearly $D$ is central and $E/D\cong H$.

It remains to split the additive sequence.  Lift the displayed generators
of $H$ to $K\subseteq E$.  If $e_i$ lifts the generator of order $n_i$,
then $n_i e_i$ is a multiple of $p^{r-1}d$.  By the choice of $r$, this
multiple equals $n_i d_i$ for some $d_i\in D$: indeed, if $p^s$ is the
$p$-part of $n_i$, then $n_iD=p^sD$ contains $p^{r-1}D$.  Replacing $e_i$
by $e_i-d_i$ makes $n_i e_i=0$.  These adjusted lifts, together with
arbitrary lifts of the free generators, define an additive section $H\to E$.
\end{proof}

The following concrete quotient remembers the central subgroup while turning
brackets into the ordinary module identity displayed at the end of the lemma.

\begin{lemma}[the central derivative module]\label{lem:central-module}
Let $D$ be a finite cyclic central ideal of a Lie ring $E$, and suppose that
$E=D\oplus H$ as an abelian group, where $H$ is finitely generated as an
abelian group.  Put
\[
 M=\mathfrak d(E)/D\mathfrak d(E),                       \tag{3}\label{eq:M}
\]
where $D\mathfrak d(E)$ is the additive subgroup spanned by the products
$dv$ with $d\in D$ and $v\in\mathfrak d(E)$.  Then $D$ embeds in $M$ as an
additive direct summand.  Moreover, $M$ is an $H$-module under
\[
 x\bar v=\overline{\widetilde x\,v},                   \tag{4}\label{eq:action}
\]
where $\widetilde x$ is any lift of $x$ to $E$.  This action annihilates
$D$, and for $e,e'\in E$, with images $x,y\in H$, one has
\[
                \overline{[e,e']}=x\bar e'-y\bar e.      \tag{5}\label{eq:derivative}
\]
\end{lemma}

\begin{proof}
Choose an additive decomposition
\[
 E=\langle d\rangle\oplus\langle e_1\rangle\oplus\cdots
   \oplus\langle e_m\rangle
\]
into cyclic groups, with $d$ generating $D$, and use the displayed order.
Replacing a pair in the wrong order by
\[
                         vu=uv+[v,u].                     \tag{6}\label{eq:collect}
\]
reduces either word length or the number of inversions, so ordered words
span $\U(E)$.  The integral PBW theorem for a direct sum of cyclic groups
says
\[
                         \operatorname{gr}\U(E)\cong\operatorname{Sym}(E).
\]
In the present coordinates this also proves uniqueness.  In a relation
between collected words, take the terms of greatest length.  Their symbols
are distinct monomials in $\operatorname{Sym}(E)$, so each coefficient is
zero modulo the greatest common divisor of the finite orders of the letters
in that word; if all those letters have infinite order, the coefficient is
zero.  The same modulus already annihilates the word in $\U(E)$.
Delete the greatest-length terms and continue downwards.

Since $d$ is central, collecting $dv$, with $v\in\mathfrak d(E)$, produces
only ordered words of length at least two containing $d$.  Conversely, every
such word belongs to $D\mathfrak d(E)$.  After quotienting, the lone word
$d$ therefore survives with its original order, and taking its
coefficient gives an additive retraction $M\to D$.

Since $D$ is central, left multiplication preserves
$D\mathfrak d(E)$, because $e(dv)=d(ev)$.  Formula \eqref{eq:action} is
independent of the lift, because two lifts differ by an element of $D$.  Also
$[\widetilde x,\widetilde y]-\widetilde{[x,y]}\in D$, so the commutator of
the two operators in \eqref{eq:action} is the operator belonging to
$[x,y]$.  Thus this is an $H$-action.  Finally,
$\widetilde x d=d\widetilde x$ vanishes in $M$.  To prove
\eqref{eq:derivative}, choose $e$ and $e'$ themselves as the lifts of $x$
and $y$; then
\[
 x\bar e'-y\bar e=\overline{ee'-e'e}=\overline{[e,e']},
\]
which proves the remaining assertions.
\end{proof}

\begin{lemma}[embedding the central extension]\label{lem:central-embedding}
Under the hypotheses of Lemma~\ref{lem:central-module}, the Lie ring $E$
embeds in
\[
                         \Hom_{\Z}(\U(H),D)\rtimes H.
\]
The image of $D$ lies in the first factor and is annihilated by the
augmentation ideal of $\U(H)$.
\end{lemma}

\begin{proof}
Put $J=\Hom_{\Z}(\U(H),D)$ and let $\varepsilon$ be the augmentation.  The
map
\[
                         j(d)(u)=\varepsilon(u)d
\]
embeds the trivial $H$-module $D$ in $J$.  By
Lemma~\ref{lem:central-module}, $D$ is an additively split submodule of $M$,
so Lemma~\ref{lem:coinduced-extension} extends $j$ to an $H$-module map
$T:M\to J$.  Define
\[
                         \Phi(e)=(T(\bar e),e+D).
\]
Equation \eqref{eq:derivative} says exactly that $\Phi$ preserves brackets.
If $\Phi(e)=0$, then $e\in D$ and $T(\bar e)=j(e)$; evaluation at $1$ gives
$e=0$.  Thus $\Phi$ is injective.  Finally, for $x\in H$,
$(xj(d))(u)=j(d)(ux)=\varepsilon(u)\varepsilon(x)d=0$, so the augmentation
ideal annihilates $j(D)$.
\end{proof}

\begin{lemma}[the augmentation Rees ring]\label{lem:rees-noetherian}
If $H$ is finitely generated and nilpotent and $I=\mathfrak d(H)$, then
$\bigoplus_{n\ge0}I^nt^n$ is Noetherian.
\end{lemma}

\begin{proof}
Choose a finite-rank free nilpotent Lie ring $F\twoheadrightarrow H$.  For a
homogeneous Hall basis $e_1,\ldots,e_s$ of respective weights
$w_1,\ldots,w_s$, weighted PBW collection gives
\[
 \mathfrak d(F)^n=
 \left\langle e_1^{a_1}\cdots e_s^{a_s}:
                   \sum_iw_i a_i\ge n\right\rangle_{\Z}.
\]
Filter the Rees ring of $F$ by the number of PBW factors.  Its associated
graded ring is commutative and is generated over $\Z$ by the classes of
$e_it^j$, where $0\le j\le w_i$: the displayed inequality says exactly that
the exponent of $t$ can be distributed among the PBW factors in this way.
The associated graded ring is therefore Noetherian by the Hilbert basis
theorem, and the usual leading-term argument makes the Rees ring of $F$
Noetherian.  The epimorphism $F\to H$ maps every augmentation power onto the
corresponding power for $H$, whose Rees ring is therefore a quotient.
\end{proof}

\begin{lemma}[separating a trivial submodule]\label{lem:rees-separation}
Let $I$ be an ideal of a ring $U$ whose Rees ring is Noetherian.  Let $V$ be
a finitely generated $U$-module and let $B\le V$ be a $U$-submodule with
$IB=0$.  Then $B\cap I^rV=0$ for some $r$.
\end{lemma}

\begin{proof}
The Rees module $\bigoplus_{n\ge0}I^nVt^n$ is finitely generated over the
Rees ring, hence is Noetherian.  Its graded submodule
\[
                    \bigoplus_{n\ge0}(B\cap I^nV)t^n
\]
has finitely many homogeneous generators.  Choose $r$ larger than all their
degrees.  In degree $r$, every generator must be multiplied by a
positive-degree coefficient, hence by an element of $I$.  Such a coefficient
kills $B$, so the degree-$r$ part is zero.
\end{proof}

\begin{lemma}[triangular realization]\label{lem:triangular}
Let $H$ act on an abelian group $V$, and let $I$ be the augmentation ideal of
$\U(H)$.  If $I^rV=0$, then $V\rtimes H$ maps to the commutator Lie ring of
a nilpotent associative ring, and this map is injective on $V$.
\end{lemma}

\begin{proof}
Write $\theta:H\to\operatorname{End}_{\Z}(V)$ for the action and send
\[
 (v,x)\longmapsto
 \begin{pmatrix}0&0\\ v&\theta(x)\end{pmatrix}
 \quad\text{in }\operatorname{End}_{\Z}(\Z\oplus V).
\]
The lower-left block of a commutator is
$\theta(x)w-\theta(y)v$, and the lower-right block is $\theta([x,y])$, so
this is a Lie homomorphism.  It is injective on $V$ because $v$ is its
lower-left block.  A product of $r+1$ displayed matrices has $r$ action maps
in its lower-left block and $r+1$ in its lower-right block.  Both vanish
because $I^rV=0$.  Thus the nonunital associative ring generated by these
matrices is nilpotent.
\end{proof}

\begin{proof}[Proof of Lemma~\ref{lem:separate}]
Put $H=K/A$, $U=\U(H)$, and $I=\mathfrak d(H)$.  The quotient $H$ is
finitely generated and nilpotent, so Lemma~\ref{lem:additive-finite} says
that its additive group is finitely generated.  Lemma~\ref{lem:finite-pushout}
places $K$ inside a Lie ring $E=D\oplus H$ with $D$ finite cyclic and
central.  Lemma~\ref{lem:central-embedding} then embeds $E$, and hence $K$,
in $\Hom_{\Z}(U,D)\rtimes H$.  The ideal $I$ annihilates the image of $A$
in the first factor.

By Lemma~\ref{lem:additive-finite}, choose finitely many additive generators
of $K$, and let $V$ be the $U$-module generated by their first coordinates.
Then $V$ is finitely generated over $U$, the embedded copy of $K$ lies in
$V\rtimes H$, and the image of $A$ lies in $V$.  Since $I$ annihilates that
image and $U/I\cong\Z$, it is a $U$-submodule of $V$.

Lemmas~\ref{lem:rees-noetherian} and~\ref{lem:rees-separation} give $r$ with
$A\cap I^rV=0$, where $A$ is identified with its image in $V$.  After
replacing $V$ by $V/I^rV$, the map from $K$ is still injective on $A$, while
$I^rV=0$.  Lemma~\ref{lem:triangular} now supplies the required nilpotent
associative ring and Lie homomorphism.
\end{proof}

\begin{lemma}[finite tail]\label{lem:finite-tail}
If $K$ is finitely generated and nilpotent of class $c$, then
$\dlt_{c+1}(K)$ is finite.
\end{lemma}

\begin{proof}
By Lemma~\ref{lem:additive-finite}, the additive group of $K$ is finitely
generated, so its torsion subgroup is finite.  The map $K\to K\otimes\Q$
respects augmentation powers.  Over $\Q$, give a basis adapted to the lower
central series its lower-central weights.  Filtered PBW then says that the
degree-one elements in the $r$th augmentation power are precisely the
elements of the $r$th lower central term.  Hence
\[
 \dlt_{c+1}(K\otimes\Q)=\gm_{c+1}(K\otimes\Q)=0.
\]
Thus $\dlt_{c+1}(K)$ lies in the kernel of rationalization, which is the
finite torsion subgroup of $K$.
\end{proof}

\begin{lemma}[zero intersection]\label{lem:zero-intersection}
If $K$ is finitely generated and nilpotent, then $\dlt_\omega(K)=0$.
\end{lemma}

\begin{proof}
First, $\dlt_\omega(K)$ is an ideal.  Indeed, if
$x\in\dlt_\omega(K)$ and $y\in K$, then for every $q\ge1$ the image of
$[x,y]=xy-yx$ in $\U(K)$ lies in
\[
 [\mathfrak d(K)^q,\mathfrak d(K)]\subseteq\mathfrak d(K)^{q+1}.
\]
Thus $[x,y]$ belongs to every dimension subring.

Suppose that $\dlt_\omega(K)$ were nonzero.  Repeatedly bracket it with $K$.
The term obtained after $i$ brackets lies in $\gm_{i+1}(K)$, so nilpotence
makes the sequence reach zero; ideality keeps every term inside
$\dlt_\omega(K)$.  Its last nonzero term is central because its bracket with
$K$ is the next, zero, term.  It lies in
$\dlt_\omega(K)\subseteq\dlt_{c+1}(K)$, where $c$ is the class of $K$, so it
is finite by Lemma~\ref{lem:finite-tail}.  It therefore contains a subgroup
$A\cong\Z/p$ for some prime $p$.

Lemma~\ref{lem:separate} gives a nilpotent associative ring $R$ and a Lie
map $\rho:K\to R^{(-)}$ which is injective on $A$.  Choose $N$ with $R^N=0$.
After adjoining an identity to $R$, the universal property extends $\rho$ to
an algebra map $\U(K)\to\Z\ltimes R$.  Since $\mathfrak d(K)$ is generated
by the image of $K$, it maps into $R$; hence $\mathfrak d(K)^N$ maps to zero.
Thus $\rho$ kills $\dlt_N(K)$, including
$A\subseteq\dlt_\omega(K)$, a contradiction.
\end{proof}

\begin{lemma}\label{lem:nilpotent}
If $K$ is finitely generated and nilpotent, then $\dlt_s(K)=0$ for some $s$.
\end{lemma}

\begin{proof}
Let $c$ be the nilpotency class.  By Lemma~\ref{lem:finite-tail},
\[
 \dlt_{c+1}(K)\supseteq\dlt_{c+2}(K)\supseteq\cdots
\]
is a descending chain in a finite group, so it stabilizes.  Its stable value
is its intersection, which is zero by Lemma~\ref{lem:zero-intersection}.
\end{proof}

\begin{proof}[Proof of Theorem~\ref{thm:main}]
For $n=1$, take $m=1$.  Now fix $n\ge2$ and put
$K=L/\gm_n(L)$.  This quotient is finitely generated, and its $n$th lower
central term is zero, so it is nilpotent.  By
Lemma~\ref{lem:nilpotent}, choose $m$ such that $\dlt_m(K)=0$.

Let $x\in\dlt_m(L)$.  The quotient map $L\to K$ induces a map of
enveloping algebras which carries $\mathfrak d(L)^m$ into
$\mathfrak d(K)^m$.  The image of $x$ therefore belongs to
$\dlt_m(K)=0$.  In other words, $x$ lies in the kernel of $L\to K$, which
is $\gm_n(L)$.  This proves \eqref{eq:main}.

If $x\in\dlt_\omega(L)$, then for each $n$ it belongs to the particular
$\dlt_m(L)$ just chosen, and hence to $\gm_n(L)$.  Thus
$\dlt_\omega(L)\subseteq\gm_\omega(L)$.  Conversely, the image of $L$ in
$\U(L)$ lies in $\mathfrak d(L)$, and
\[
 [\mathfrak d(L)^i,\mathfrak d(L)^j]
        \subseteq\mathfrak d(L)^{i+j}.
\]
Induction on $n$ therefore gives $\gm_n(L)\subseteq\dlt_n(L)$.  Taking
intersections gives the reverse inclusion and completes the proof.
\end{proof}

\end{document}
